% Look-ahead report
% Tables under tables/ are generated by tests/make_tables_lookahead.py from
% results/lookahead.csv, so `make bench-lookahead` refreshes them

\documentclass[11pt,a4paper]{article}
% Shared preamble for the NAIL094 reports
% Used by every task via % Shared preamble for the NAIL094 reports
% Used by every task via % Shared preamble for the NAIL094 reports
% Used by every task via \input{../../shared/preamble}

\usepackage[T1]{fontenc}
\usepackage[utf8]{inputenc}
\usepackage{lmodern}
\usepackage[margin=2.5cm]{geometry}

\usepackage{amsmath}
\usepackage{amssymb}
\usepackage{booktabs}
\usepackage{float}
\usepackage{graphicx}
\usepackage{listings}
\usepackage{xcolor}
\usepackage{caption}

\usepackage[hidelinks]{hyperref}

\setlength{\parskip}{0.5em}
\setlength{\parindent}{0pt}

% Title fields, set these in the report before \begin{document}
\makeatletter
\newcommand{\reportcourse}[1]{\def\@reportcourse{#1}}
\newcommand{\reporttask}[1]{\def\@reporttask{#1}}
\newcommand{\reportauthor}[1]{\def\@reportauthor{#1}}

\newcommand{\makereporttitle}{%
  \begin{center}
    {\LARGE\bfseries \@reportcourse \par}
    \vspace{0.8em}
    {\Large \@reporttask \par}
    \vspace{0.6em}
    {\normalsize \@reportauthor \par}
  \end{center}
  \vspace{1em}
  \hrule
  \vspace{1.5em}
}
\makeatother

% Inline code
\newcommand{\code}[1]{\texttt{#1}}

% Code listings
\lstdefinestyle{reportcode}{
  basicstyle=\ttfamily\small,
  keywordstyle=\color{blue!60!black},
  commentstyle=\color{gray},
  stringstyle=\color{green!40!black},
  breaklines=true,
  showstringspaces=false,
  frame=single,
  rulecolor=\color{gray!40},
  numbers=none,
}
\lstset{style=reportcode}


\usepackage[T1]{fontenc}
\usepackage[utf8]{inputenc}
\usepackage{lmodern}
\usepackage[margin=2.5cm]{geometry}

\usepackage{amsmath}
\usepackage{amssymb}
\usepackage{booktabs}
\usepackage{float}
\usepackage{graphicx}
\usepackage{listings}
\usepackage{xcolor}
\usepackage{caption}

\usepackage[hidelinks]{hyperref}

\setlength{\parskip}{0.5em}
\setlength{\parindent}{0pt}

% Title fields, set these in the report before \begin{document}
\makeatletter
\newcommand{\reportcourse}[1]{\def\@reportcourse{#1}}
\newcommand{\reporttask}[1]{\def\@reporttask{#1}}
\newcommand{\reportauthor}[1]{\def\@reportauthor{#1}}

\newcommand{\makereporttitle}{%
  \begin{center}
    {\LARGE\bfseries \@reportcourse \par}
    \vspace{0.8em}
    {\Large \@reporttask \par}
    \vspace{0.6em}
    {\normalsize \@reportauthor \par}
  \end{center}
  \vspace{1em}
  \hrule
  \vspace{1.5em}
}
\makeatother

% Inline code
\newcommand{\code}[1]{\texttt{#1}}

% Code listings
\lstdefinestyle{reportcode}{
  basicstyle=\ttfamily\small,
  keywordstyle=\color{blue!60!black},
  commentstyle=\color{gray},
  stringstyle=\color{green!40!black},
  breaklines=true,
  showstringspaces=false,
  frame=single,
  rulecolor=\color{gray!40},
  numbers=none,
}
\lstset{style=reportcode}


\usepackage[T1]{fontenc}
\usepackage[utf8]{inputenc}
\usepackage{lmodern}
\usepackage[margin=2.5cm]{geometry}

\usepackage{amsmath}
\usepackage{amssymb}
\usepackage{booktabs}
\usepackage{float}
\usepackage{graphicx}
\usepackage{listings}
\usepackage{xcolor}
\usepackage{caption}

\usepackage[hidelinks]{hyperref}

\setlength{\parskip}{0.5em}
\setlength{\parindent}{0pt}

% Title fields, set these in the report before \begin{document}
\makeatletter
\newcommand{\reportcourse}[1]{\def\@reportcourse{#1}}
\newcommand{\reporttask}[1]{\def\@reporttask{#1}}
\newcommand{\reportauthor}[1]{\def\@reportauthor{#1}}

\newcommand{\makereporttitle}{%
  \begin{center}
    {\LARGE\bfseries \@reportcourse \par}
    \vspace{0.8em}
    {\Large \@reporttask \par}
    \vspace{0.6em}
    {\normalsize \@reportauthor \par}
  \end{center}
  \vspace{1em}
  \hrule
  \vspace{1.5em}
}
\makeatother

% Inline code
\newcommand{\code}[1]{\texttt{#1}}

% Code listings
\lstdefinestyle{reportcode}{
  basicstyle=\ttfamily\small,
  keywordstyle=\color{blue!60!black},
  commentstyle=\color{gray},
  stringstyle=\color{green!40!black},
  breaklines=true,
  showstringspaces=false,
  frame=single,
  rulecolor=\color{gray!40},
  numbers=none,
}
\lstset{style=reportcode}


\reportcourse{NAIL094 --- Decision Procedures and SAT/SMT Solvers}
\reporttask{Look-Ahead Solver}
\reportauthor{Zerina Ahmetovic}

\begin{document}
\makereporttitle

\section{Task overview}

The branching rule of the DPLL solver~\cite{dpllreport} is replaced by a
look-ahead procedure. Three difference heuristics are implemented and compared,
two additional reasoning techniques are added, and both rely on an eager
clause-length data structure. Look-ahead is switched on with
\texttt{dpll --lookahead [--heuristic=crh|wbh|bsh] [input]}.

This is an addition to the previous DPLL algorithm~\cite{dpllreport}, not a
separate solver: run without \texttt{--lookahead} still branches in index
order and behaves exactly as that report describes. The two sets of numbers below therefore come
from one program, not two, and the branching rule is the only thing that
differs between them.

The benchmarks are the same non-random SATLIB~\cite{satlib} families used
before, pigeonhole~\cite{phole} and All Interval Series~\cite{ais}, with one
uniform random 3-SAT~\cite{uf} instance for contrast.

\section{Implementation}

At every node each free variable is tentatively assigned in both polarities and
propagated. A conflict in one polarity makes the opposite literal a failed
literal, which is forced; a conflict in both makes the node unsatisfiable.
Otherwise the reduction is measured and undone.

\textbf{Difference heuristics.} All three follow the definitions of Heule and
van Maaren~\cite{handbook}. Let $\#_k(l)$ be the number of occurrences of
literal $l$ in clauses of size $k$ in the input formula, and
$w(l) = \sum_{k \ge 2} \gamma_k \cdot \#_k(l)$.

\begin{itemize}
\item CRH weighs each newly created clause by its size, using Kullmann's
      constants $\gamma_2 = 1$, $\gamma_3 = 0.2$, $\gamma_4 = 0.05$,
      $\gamma_5 = 0.01$, $\gamma_6 = 0.003$, and above that the regression
      $\gamma_k = 20.4514 \cdot 0.218673^{\,k}$ these were fitted to.
      Both families reach into that range --- pigeonhole has one clause per
      pigeon of size $n$, All Interval Series $n$ of size $n$ and $n-1$ of
      size $n-1$ --- but the weights there are vanishing: a clause of size
      $8$ counts about a nine-thousandth of a new binary clause and one of
      size $12$ about a four-millionth, so CRH is in effect blind to them.
\item WBH sums $w(\lnot a) + w(\lnot b)$ over the newly created binary clauses
      $(a \vee b)$, with $\gamma_k = 5^{3-k}$.
\item BSH multiplies instead of adding, $w(\lnot a) \cdot w(\lnot b)$, with
      $\gamma_k = 2^{3-k}$.
\end{itemize}

The exponents used here are the ones given in the task
assignment~\cite{taskpage}: $\gamma_k = 5^{3-k}$ for WBH and
$\gamma_k = 2^{3-k}$ for BSH. The weights $w(l)$ are computed once from the
input formula and not recomputed as the search descends, which satz does; that
is a simplification of the chapter's scheme rather than a faithful copy of
it.

The two directions are combined by the MixDiff function of the same chapter,
$1024 \cdot L \cdot R + L + R$, where $L$ and $R$ are the values for the two
polarities. The product rewards variables that constrain both ways and the sum
only breaks ties. The more reduced polarity is explored first, being the
smaller subproblem.

\textbf{Additional reasoning.} Two techniques are implemented. \emph{Local
learning}: when the look-ahead on $x$ derives $y$, the binary clause
$\lnot x \vee y$ holds under the current partial assignment and is added; such
clauses are propagated like any other and are dropped again when the trail
shrinks below the point at which they were added. \emph{Autarky reasoning}: if
a look-ahead creates no new clause, every clause it touched came out satisfied,
so the assignment is an autarky and can be committed without branching.

\textbf{Eager data structure.} Propagation keeps, per clause, the number of its
literals that are true and the number that are not false. The second is the
clause's current length, which is exactly what the heuristics read. Watched
literals deliberately do not track it, so the look-ahead solver uses this third
engine instead; the cost is that backtracking has to undo the counters, where
the other two engines undo nothing at all.

The engine can also be selected on its own with \texttt{--counters}, which
keeps the counters but branches in index order. Although this way would not be the
natural use of the solver, it is kept as it is the only configuration in which
the data structure varies and the search does not, so it is what makes the
cost of the counters measurable.

\section{Experimental evaluation}

The measurements were taken with g++ 13.3.0 under WSL2 (Ubuntu 24.04) at
\texttt{-std=c++20 -O2}. Each instance was run three times and the fastest
kept, the timer covers the search only. All configurations agree on every
instance, and reported models were checked against the formula.

\begin{table}[H]
\centering
\small
\setlength{\tabcolsep}{5pt}
\begin{tabular}{lrlrrrrrr}
\toprule
instance & vars & result & plain DPLL & counters & CRH & WBH & BSH & WBH bare \\
\midrule
hole6 & 42 & UNSAT & 6\,490 & 6\,490 & 478 & 478 & 478 & 1\,038 \\
hole7 & 56 & UNSAT & 65\,560 & 65\,560 & 3\,358 & 3\,358 & 3\,358 & 10\,498 \\
hole8 & 72 & UNSAT & 756\,686 & 756\,686 & 26\,878 & 26\,878 & 26\,878 & 121\,182 \\
ais6 & 61 & SAT & 57 & 57 & 4 & 3 & 6 & 3 \\
ais8 & 113 & SAT & 550 & 550 & 6 & 5 & 5 & 6 \\
ais10 & 181 & SAT & 8\,973 & 8\,973 & 13 & 10 & 12 & 13 \\
ais12 & 265 & SAT & 192\,780 & 192\,780 & 29 & 6 & 6 & 6 \\
uf50-01 & 50 & SAT & 1\,016 & 1\,016 & 9 & 9 & 4 & 13 \\
\bottomrule
\end{tabular}
\caption{Number of decisions, that is the size of the search tree. The counters column is the eager engine under plain branching, a control rather than a solver variant; its tree is identical to plain DPLL by construction. WBH bare is WBH with local learning and autarky reasoning switched off.}
\label{tab:decisions}
\end{table}


Table~\ref{tab:decisions} shows the search tree collapsing. Against plain DPLL
the reduction is a factor of $14$ to $28$ on pigeonhole and far larger on All
Interval Series, where \texttt{ais12} falls from $192\,780$ decisions to $6$
under WBH.

That gap is a difference in kind rather than degree, and it follows from the
verdicts. Pigeonhole is unsatisfiable, so every branch has to be refuted and
none can be skipped; a better branching variable shortens the refutation but
cannot avoid doing one, which bounds how much any heuristic can win. All
Interval Series is satisfiable, so a heuristic that keeps choosing productive
variables can descend almost directly to a model, and six decisions on
\texttt{ais12} is what that looks like.

The counters column repeats plain DPLL on every row. The same index-order branching over the same formula must produce the
same tree, so any difference in time between those two columns is the data
structure and nothing else.

The three heuristics separate only where the instance gives them room to. On
pigeonhole all three agree exactly, at $478$, $3\,358$ and $26\,878$ decisions:
the formula consists of binary clauses and one long clause per pigeon, the long
clauses carry almost no weight at any of the three weightings, and the
remaining binary structure is symmetric enough that every candidate ties. The
difference appears on All Interval Series, where WBH and BSH find
\texttt{ais12} in $6$ decisions against CRH's $29$.

The single random instance in Table~\ref{tab:decisions} settles nothing either
way, so the three were also run over $300$ instances of the SATLIB uniform
random 3-SAT set, all $50$ variables and $218$ clauses. Look-ahead cuts the
mean tree there from $418.2$ decisions to about four, but it does not separate
the heuristics: the means are $4.11$, $4.07$ and $4.30$ for CRH, WBH and BSH,
and every pairwise difference taken over the same instances has a $95\%$
confidence interval straddling zero. The ordering is not even stable --- over
the first hundred instances BSH led, over three hundred it trails. The choice
between the three shows up on structured instances and not on random ones.

The last column shows what the additional reasoning is worth (with local
learning and autarky reasoning switched off), \texttt{hole8} needs $121\,182$
decisions instead of $26\,878$, more than four times as many.

\begin{table}[H]
\centering
\small
\setlength{\tabcolsep}{5pt}
\begin{tabular}{lrlrrrrrr}
\toprule
instance & vars & result & plain DPLL & counters & CRH & WBH & BSH & WBH bare \\
\midrule
hole6 & 42 & UNSAT & 0.0008 & 0.0016 & 0.0056 & 0.0061 & 0.0059 & 0.0066 \\
hole7 & 56 & UNSAT & 0.0095 & 0.0122 & 0.0470 & 0.0518 & 0.0529 & 0.0751 \\
hole8 & 72 & UNSAT & 0.1192 & 0.1611 & 0.4931 & 0.5440 & 0.5361 & 1.0118 \\
ais6 & 61 & SAT & 0.0001 & 0.0001 & 0.0007 & 0.0008 & 0.0009 & 0.0005 \\
ais8 & 113 & SAT & 0.0009 & 0.0010 & 0.0036 & 0.0027 & 0.0034 & 0.0022 \\
ais10 & 181 & SAT & 0.0144 & 0.0157 & 0.0253 & 0.0139 & 0.0175 & 0.0088 \\
ais12 & 265 & SAT & 0.3563 & 0.3707 & 0.1835 & 0.0204 & 0.0205 & 0.0115 \\
uf50-01 & 50 & SAT & 0.0008 & 0.0011 & 0.0009 & 0.0008 & 0.0006 & 0.0008 \\
\bottomrule
\end{tabular}
\caption{Running time in seconds for the same runs, fastest of three, search only.}
\label{tab:times}
\end{table}


Table~\ref{tab:times} shows that a smaller tree is not always a faster solve.
On pigeonhole look-ahead loses: \texttt{hole8} takes $0.4931$\,s under CRH
against plain DPLL's $0.1192$\,s, because probing every free variable at every
node costs more than the $28$-fold saving in nodes is worth. On All Interval
Series it wins decisively, \texttt{ais12} dropping from $0.3563$\,s to
$0.0204$\,s under WBH. Here the decision counts carry straight over into time:
CRH needs $29$ decisions to the others' $6$ and is correspondingly slower at
$0.1835$\,s. Two instances, \texttt{ais6} and \texttt{uf50-01}, run in under
two milliseconds in every configuration, so their ratios are not worth reading.

Because the counters column shares plain DPLL's tree, the slowdown on 
pigeonhole splits into two factors that multiply.
On \texttt{hole8} the eager engine alone costs $0.1611$\,s against
$0.1192$\,s, a factor of $1.35$, and look-ahead on top of it costs
$0.4931$\,s against $0.1611$\,s, a factor of $3.06$; together they give the
$4.14$ observed. In time, look-ahead adds $0.374$\,s to the plain solve, of
which the counters are $0.042$\,s and the probing $0.332$\,s --- so the
bookkeeping is $11\%$ of the excess and the probing the remaining $89\%$. On All Interval Series
the engine is close to free, being a factor of $1.04$ on \texttt{ais12}, and
look-ahead still gains a factor of $18$ on top of it, $0.3707$\,s against
$0.0204$\,s.

The additional reasoning follows the same split. On pigeonhole it pays for
itself, cutting \texttt{hole8} from $1.0118$\,s to $0.5440$\,s. On All Interval
Series the tree is already down to single-digit decisions, so there is nothing
left for it to save and maintaining the learned clauses is a net cost:
\texttt{ais12} is faster without it, at $0.0115$\,s against $0.0204$\,s.

\subsection*{Code}

\begin{sloppypar}
The full source is in the accompanying repository~\cite{repo}.
\texttt{src/look\_ahead.hpp} holds the look-ahead procedure and the three
heuristics, \texttt{src/eager\_counters.hpp} the counter-based propagation and
the learned clause store. Both include \texttt{src/dpll.hpp} and nothing in
\texttt{dpll.hpp} refers back to them.
The \texttt{--counters} flag selecting the engine under plain branching is in
\texttt{src/dpll.cpp}.
\texttt{tests/check\_lookahead.py} asserts that every configuration agrees and
validates the reported models, \texttt{tests/fuzz.py} checks each heuristic
against exhaustive search on random instances, and
\texttt{tests/bench\_lookahead.py} produces the measurements.
\end{sloppypar}

\subsection*{Summary}

Look-ahead buys a much smaller search tree at a large cost per node, and
whether that is worth it follows the verdict. It wins by more than an order of
magnitude on the satisfiable All Interval Series instances, where a good
heuristic descends almost straight to a model; it loses on the unsatisfiable
pigeonhole ones, where every branch must be refuted, the tree shrinks by only a
factor of $28$ and the solve still takes four times longer. Running the eager
engine under plain branching splits that loss: on \texttt{hole8} the counters
account for $11\%$ of it and the probing for the rest.

The choice between the three heuristics matters far less than either. They are
identical on pigeonhole, statistically indistinguishable over $300$ random
instances, and separate only on All Interval Series, while the additional
reasoning makes a larger difference than any of them. No configuration measured
here is best everywhere, and the ones that win do so for reasons visible in the
formula rather than by being better in general.

\begin{thebibliography}{9}

\bibitem{dpllreport}
Z.~Ahmetovic.
\newblock DPLL experimental evaluation, an earlier report of this series.
\newblock Included in the accompanying repository.

\bibitem{repo}
Z.~Ahmetovic.
\newblock Source repository for the NAIL094 programming tasks.
\newblock \url{https://github.com/inferno73/SAT-solver-NAIL094}

\bibitem{taskpage}
P.~Ku\v{c}era.
\newblock NAIL094 programming task: Look-Ahead Solver.
\newblock \url{https://ktiml.mff.cuni.cz/~kucerap/satsmt/practical/task_look_ahead.php}

\bibitem{handbook}
M.~Heule and H.~van Maaren.
\newblock Look-Ahead Based SAT Solvers.
\newblock In \emph{Handbook of Satisfiability}, chapter 5, pp. 155--184.
\newblock Frontiers in Artificial Intelligence and Applications, vol. 185,
IOS Press, 2009.
\newblock \url{https://www.cs.utexas.edu/~marijn/publications/p01c05_lah.pdf}

\bibitem{satlib}
H.~H.~Hoos and T.~St\"utzle.
\newblock SATLIB: An Online Resource for Research on SAT.
\newblock \url{https://www.cs.ubc.ca/~hoos/SATLIB/benchm.html}

\bibitem{ais}
SATLIB benchmark problems, All Interval Series.
\newblock \url{https://www.cs.ubc.ca/~hoos/SATLIB/Benchmarks/SAT/AIS/descr.html}

\bibitem{phole}
SATLIB benchmark problems, DIMACS pigeon hole.
\newblock \url{https://www.cs.ubc.ca/~hoos/SATLIB/Benchmarks/SAT/DIMACS/PHOLE/descr.html}

\bibitem{uf}
SATLIB benchmark problems, Uniform Random-3-SAT.
\newblock \url{https://www.cs.ubc.ca/~hoos/SATLIB/Benchmarks/SAT/RND3SAT/descr.html}

\end{thebibliography}

\end{document}
